% Generated by tools/edn_latex.py from reports/edn.md.
% Do not edit: change reports/edn.md and run `make edn`.
\ednentry{
  id = {EDN-10},
  title = {The availability target is recovery within 30 seconds},
  date = {2026-09-22},
  milestone = {M6: Monitoring},
  topic = {Resource Monitoring, Requirements},
  participants = {@lukas2510},
  decision = {NFR-05 no longer promises 99\,\% uptime.
The API recovers by itself within 30 seconds of a crash or a reboot, which a rehearsed crash verifies.
It has no unplanned downtime during a graded presentation and its warm-up.
Outside those windows, uptime is monitored with Better Uptime and Grafana and reported, not promised.},
  rationale = {\emph{Alternatives considered:}
\begin{itemize}[leftmargin=*, topsep=2pt]
\item Keep 99\,\% uptime: nothing backs it on one VM without redundancy, SLA, or on-call, and its measurement window was unbounded.
\item Lower the number to 95\,\%: as unverifiable as 99\,\%.
\end{itemize}
\par
\emph{Why:} M6 grades whether monitoring gives actionable insight, not whether a percentage was hit.
The new requirement separates what we control, recovery and the presentation days, from what we can only observe.
The recovery clause is testable today.},
  ai = {Information seeking, Alternative generation, Alternative assessment, Recommendation},
  response = {Accepted},
  assessment = {AI had flagged that Prometheus's \texttt{up} metric cannot see the host going down and that Better Uptime, a course tool, was missing from the docs.
Asked which target is feasible, it proposed the three-part requirement.
Lukas confirmed it.},
  aievidence = {Claude Code session of 2026-09-22, review of PR \#19; Lukas asked: ``what NFR that goes in the uptime direction is feasible with our system and constraints?''.},
  evidence = {\href{https://github.com/mlops-2627q1-mds-upc/MLOPS_Recommenditos/blob/main/docs/docs/requirements.md}{Requirements} NFR-05, \href{https://github.com/mlops-2627q1-mds-upc/MLOPS_Recommenditos/blob/main/docs/docs/specification.md}{specification} NFR-05, \href{https://github.com/mlops-2627q1-mds-upc/MLOPS_Recommenditos/pull/19}{PR \#19}.},
}
